\chapter{Numerical Differentiation and Integration}

\section{Numerical Differentiation}

Taylor expansions about $x$ give
\[
 f(x+h)=f(x)+hf'(x)+\frac{h^2}{2!}f''(x)
 +\frac{h^3}{3!}f^{(3)}(x)+\cdots,
\]
\[
 f(x-h)=f(x)-hf'(x)+\frac{h^2}{2!}f''(x)
 -\frac{h^3}{3!}f^{(3)}(x)+\cdots.
\]
By adding, subtracting or rearranging these series, derivatives can be approximated by finite differences.

\begin{definition}[Finite-difference formulas]
For a uniform step $h$,
\[
 f'(x)=\frac{f(x+h)-f(x)}h+\bigO(h)
 \qquad\text{(forward difference)},
\]
\[
 f'(x)=\frac{f(x)-f(x-h)}h+\bigO(h)
 \qquad\text{(backward difference)},
\]
\[
 f'(x)=\frac{f(x+h)-f(x-h)}{2h}+\bigO(h^2)
 \qquad\text{(central difference)},
\]
and
\[
 f''(x)=\frac{f(x+h)-2f(x)+f(x-h)}{h^2}+\bigO(h^2).
\]
\end{definition}

The order term describes the truncation error as $h\to0$.  A central difference is normally more accurate than a one-sided difference with the same step size, provided data are available on both sides of the point.

\begin{example}
Given
\[
\begin{array}{c|cccc}
\toprule
x&1.00&1.01&1.02&1.03\\
\midrule
f(x)&5.00&6.01&7.04&8.09\\
\bottomrule
\end{array}
\]
use $h=0.01$ to:
\begin{enumerate}[label=(\alph*),leftmargin=2.7em]
  \item estimate $f'(1.02)$ by forward, backward and central differences;
  \item estimate $f''(1.02)$ by a central difference;
  \item calculate the signed percentage relative errors in part (a) if the exact derivative is $104$.
\end{enumerate}
\end{example}

\begin{solution}[7.0cm]
The forward estimate is
\[
 \frac{f(1.03)-f(1.02)}{0.01}
 =\frac{8.09-7.04}{0.01}=105.
\]
The backward estimate is
\[
 \frac{f(1.02)-f(1.01)}{0.01}
 =\frac{7.04-6.01}{0.01}=103.
\]
The central estimate is
\[
 \frac{f(1.03)-f(1.01)}{0.02}
 =\frac{8.09-6.01}{0.02}=104.
\]
For the second derivative,
\[
 f''(1.02)\approx
 \frac{8.09-2(7.04)+6.01}{(0.01)^2}=200.
\]
Using
\[
 \frac{104-\widehat f'}{104}\times100\%,
\]
the signed errors are
\[
 -0.9615\%,
 \qquad 0.9615\%,
 \qquad 0\%.
\]
\end{solution}

\subsection{Richardson Extrapolation}

Suppose an approximation has an even-power error expansion
\[
 D(h)=D+c_1h^2+c_2h^4+c_3h^6+\cdots.
\]
Combining estimates at $h$ and $h/2$ eliminates the leading $h^2$ term.  Define
\[
 D_{n,0}=D\left(\frac{h}{2^n}\right).
\]
The Richardson table is generated by
\[
 D_{n,m}
 =D_{n,m-1}
 +\frac{D_{n,m-1}-D_{n-1,m-1}}{4^m-1}.
\]
The entry $D_{n,m}$ has formal error order $\bigO(h^{2(m+1)})$.

\begin{example}
Let $f(x)=xe^x$.  Using the central difference and Richardson extrapolation, approximate $f'(2)$ with orders $\bigO(h^2)$, $\bigO(h^4)$ and $\bigO(h^6)$, starting from $h=0.2$.
\end{example}

\begin{solution}[7.2cm]
The central-difference values are
\[
\begin{aligned}
D_{0,0}&=D(0.2)=22.414160657,\\
D_{1,0}&=D(0.1)=22.228786880,\\
D_{2,0}&=D(0.05)=22.182564858.
\end{aligned}
\]
First extrapolation gives
\[
 D_{1,1}=D_{1,0}+\frac{D_{1,0}-D_{0,0}}3
 =22.166995621,
\]
\[
 D_{2,1}=D_{2,0}+\frac{D_{2,0}-D_{1,0}}3
 =22.167157517.
\]
The next extrapolation is
\[
 D_{2,2}=D_{2,1}+\frac{D_{2,1}-D_{1,1}}{15}
 =22.167168310.
\]
Thus the requested approximations are
\[
 22.414161,
 \qquad22.166996,
 \qquad22.167169.
\]
The exact derivative is $f'(2)=3e^2\approx22.16716830$.
\end{solution}

\section{Numerical Integration}

The numerical integration problem is to approximate
\[
 I=\int_a^b f(x)\dd x
\]
when an elementary antiderivative is unavailable or when $f$ is known only through data.

\subsection{Composite Trapezoidal Rule}

On one interval,
\[
 \int_a^b f(x)\dd x
 \approx\frac{b-a}{2}\bigl[f(a)+f(b)\bigr].
\]
For a uniform partition
\[
 x_i=a+ih,
 \qquad h=\frac{b-a}{n},
\]
the composite trapezoidal rule is
\[
 T_n=\frac h2
 \left[f(x_0)+2\sum_{i=1}^{n-1}f(x_i)+f(x_n)\right].
\]

\begin{theorem}[Composite trapezoidal error bound]
If $f''$ is continuous and $|f''(x)|\le M$ on $[a,b]$, then
\[
 |I-T_n|
 \le\frac{(b-a)^3}{12n^2}M.
\]
\end{theorem}

\begin{example}
Use the composite trapezoidal rule with $n=4$ to estimate
\[
 \int_1^2x^2\dd x.
\]
\end{example}

\begin{solution}[4.0cm]
Here $h=(2-1)/4=1/4$ and the nodes are
$1,1.25,1.5,1.75,2$.  Hence
\[
\begin{aligned}
T_4&=\frac{1/4}{2}
\left[1^2+2(1.25^2+1.5^2+1.75^2)+2^2\right]\\
&=\frac{75}{32}=2.34375.
\end{aligned}
\]
The exact integral is $7/3\approx2.33333$.
\end{solution}

\begin{example}
How many subintervals are sufficient to approximate
\[
 \int_1^2\frac1x\dd x
\]
by the composite trapezoidal rule with error less than $10^{-4}$?
\end{example}

\begin{solution}[3.5cm]
For $f(x)=1/x$,
\[
 f''(x)=\frac2{x^3},
\]
so $M=2$ on $[1,2]$.  The bound requires
\[
 \frac{(2-1)^3}{12n^2}(2)<10^{-4},
\]
or
\[
 n^2>\frac1{6\times10^{-4}}=1666.67.
\]
Thus $n>40.824$, and the smallest integer choice is
\[
 n=41.
\]
\end{solution}

\subsection{Composite Simpson Rule}

On two subintervals of width $h=(b-a)/2$, Simpson's rule is
\[
 \int_a^b f(x)\dd x
 \approx\frac h3\bigl[f(x_0)+4f(x_1)+f(x_2)\bigr].
\]
For an even number $n$ of uniform subintervals,
\[
 S_n=\frac h3\left[
 f(x_0)
 +4\sum_{\substack{i=1\\i\text{ odd}}}^{n-1}f(x_i)
 +2\sum_{\substack{i=2\\i\text{ even}}}^{n-2}f(x_i)
 +f(x_n)
 \right].
\]

\begin{theorem}[Composite Simpson error bound]
If $f^{(4)}$ is continuous and $|f^{(4)}(x)|\le M$ on $[a,b]$, then
\[
 |I-S_n|
 \le\frac{(b-a)^5}{180n^4}M,
\]
where $n$ is even.
\end{theorem}

\begin{example}
Approximate
\[
 \int_0^1 4x^3\dd x
\]
by composite Simpson's rule with $n=4$.
\end{example}

\begin{solution}[3.5cm]
Simpson's rule is exact for every polynomial of degree at most three.  Therefore
\[
 S_4=\int_0^1 4x^3\dd x=[x^4]_0^1=1.
\]
Direct substitution into the composite formula gives the same result.
\end{solution}

\begin{example}
Use composite Simpson's rule with $n=10$ to approximate
\[
 \int_0^1 e^{x^2}\dd x
\]
and estimate the error using the theoretical bound.
\end{example}

\begin{solution}[5.5cm]
With $h=0.1$, the composite Simpson sum gives
\[
 S_{10}\approx1.462681.
\]
For $f(x)=e^{x^2}$,
\[
 f^{(4)}(x)=(16x^4+48x^2+12)e^{x^2}.
\]
This is increasing on $[0,1]$, so
\[
 M=f^{(4)}(1)=76e.
\]
Hence
\[
 |E_S|
 \le\frac{76e}{180(10)^4}
 \approx1.15\times10^{-4}.
\]
Thus the approximation and error bound are
\[
 1.462681,
 \qquad |E_S|\le0.000115.
\]
\end{solution}

\subsection{Newton--Cotes Rules}

Newton--Cotes formulas integrate polynomial interpolants at equally spaced nodes.  Common closed rules include
\[
 \int_{x_0}^{x_1}f(x)\dd x
 =\frac h2(f_0+f_1)-\frac{h^3}{12}f''(\xi)
 \qquad\text{(trapezoidal)},
\]
\[
 \int_{x_0}^{x_2}f(x)\dd x
 =\frac h3(f_0+4f_1+f_2)-\frac{h^5}{90}f^{(4)}(\xi)
 \qquad\text{(Simpson $1/3$)},
\]
\[
 \int_{x_0}^{x_3}f(x)\dd x
 =\frac{3h}{8}(f_0+3f_1+3f_2+f_3)
 -\frac{3h^5}{80}f^{(4)}(\xi)
 \qquad\text{(Simpson $3/8$)}.
\]
An open rule omits the endpoints.  The midpoint rule, for example, is
\[
 \int_a^b f(x)\dd x
 =(b-a)f\left(\frac{a+b}{2}\right)
 +\frac{(b-a)^3}{24}f''(\xi).
\]
Open rules are useful when the integrand is undefined or poorly behaved at an endpoint.

\section{Romberg Integration}

Romberg integration applies Richardson extrapolation to a sequence of composite trapezoidal approximations.  Define
\[
 R_{n,0}=T_{2^n}.
\]
Then
\[
 R_{n,m}=R_{n,m-1}
 +\frac{R_{n,m-1}-R_{n-1,m-1}}{4^m-1}.
\]
The triangular array successively removes the $h^2,h^4,\ldots$ error terms.

\begin{example}
Use Romberg integration through $R_{3,3}$ to approximate
\[
 \int_0^1\frac{1}{1+x^2}\dd x.
\]
\end{example}

\begin{solution}[7.0cm]
The Romberg table is
\[
\begin{array}{c|cccc}
\toprule
n&R_{n,0}&R_{n,1}&R_{n,2}&R_{n,3}\\
\midrule
0&0.750000000&&&\\
1&0.775000000&0.783333333&&\\
2&0.782794118&0.785392157&0.785529412&\\
3&0.784747124&0.785398126&0.785398524&0.785396446\\
\bottomrule
\end{array}
\]
Therefore
\[
 R_{3,3}\approx0.785396.
\]
The exact value is $\arctan(1)-\arctan(0)=\pi/4\approx0.785398$.
\end{solution}

\section{Gaussian Quadrature}

Newton--Cotes rules prescribe equally spaced nodes.  Gaussian quadrature instead chooses both the nodes and weights to maximise the degree of exactness.  An $N$-point Gauss--Legendre rule is exact for polynomials of degree at most $2N-1$.

\begin{theorem}[Two-point Gauss--Legendre rule]
For a sufficiently smooth function on $[-1,1]$,
\[
 \int_{-1}^{1}f(x)\dd x
 \approx
 f\left(-\frac1{\sqrt3}\right)
 +f\left(\frac1{\sqrt3}\right).
\]
\end{theorem}

\begin{example}
Use two-point Gauss--Legendre quadrature to estimate
\[
 \int_{-1}^{1}\frac1{x+2}\dd x.
\]
\end{example}

\begin{solution}[3.5cm]
The approximation is
\[
 G_2
 =\frac1{2-1/\sqrt3}
 +\frac1{2+1/\sqrt3}
 =\frac{12}{11}
 \approx1.09091.
\]
The exact value is $\ln3\approx1.09861$.
\end{solution}

For an arbitrary interval $[a,b]$, use
\[
 x=\frac{b-a}{2}t+\frac{a+b}{2},
 \qquad
 \dd x=\frac{b-a}{2}\dd t.
\]
Thus
\[
 \int_a^b f(x)\dd x
 =\frac{b-a}{2}
 \int_{-1}^{1}
 f\left(\frac{b-a}{2}t+\frac{a+b}{2}\right)\dd t.
\]

\begin{example}
Estimate
\[
 \int_{0.1}^{1.3}5xe^{-2x}\dd x
\]
using two-point Gauss--Legendre quadrature.
\end{example}

\begin{solution}[4.5cm]
Here
\[
 x=0.6t+0.7,
 \qquad \dd x=0.6\dd t.
\]
Therefore
\[
 I\approx0.6\left[
 f\left(0.7-\frac{0.6}{\sqrt3}\right)
 +f\left(0.7+\frac{0.6}{\sqrt3}\right)
 \right],
\]
where $f(x)=5xe^{-2x}$.  Substitution gives
\[
 I\approx0.91018.
\]
\end{solution}

For reference, the first Gauss--Legendre rules on $[-1,1]$ are
\[
\begin{array}{c|c|c}
\toprule
N&\text{nodes}&\text{weights}\\
\midrule
2&\pm1/\sqrt3&1,1\\[1mm]
3&0,\ \pm\sqrt{3/5}&8/9,\ 5/9,\ 5/9\\
\bottomrule
\end{array}
\]
Higher-order rules are usually taken from a standard table or generated computationally.

Gaussian quadrature also extends to weighted integrals
\[
 \int_a^b w(x)f(x)\dd x.
\]
Important families include Gauss--Chebyshev, Gauss--Laguerre and Gauss--Hermite rules.

\begin{example}
Use the two-point Gauss--Chebyshev rule to evaluate
\[
 \int_{-1}^{1}\frac{x^2}{\sqrt{1-x^2}}\dd x.
\]
\end{example}

\begin{solution}[3.2cm]
For the weight $(1-x^2)^{-1/2}$, the two nodes are
$\pm1/\sqrt2$ and both weights are $\pi/2$.  Hence
\[
 I\approx\frac\pi2
 \left[\left(-\frac1{\sqrt2}\right)^2
 +\left(\frac1{\sqrt2}\right)^2\right]
 =\frac\pi2.
\]
The rule is exact here because $x^2$ is a polynomial of sufficiently low degree.
\end{solution}

\section{Adaptive Quadrature}

A fixed uniform mesh may waste function evaluations in smooth regions and use too few points where the integrand changes rapidly.  Adaptive quadrature subdivides only those intervals whose estimated local error is too large.

For Simpson's rule on $[a,b]$, let $S^{(1)}=S(a,b)$ and let
\[
 S^{(2)}=S\left(a,\frac{a+b}{2}\right)
 +S\left(\frac{a+b}{2},b\right).
\]
Because Simpson's error scales like $h^5$, halving the interval suggests
\[
 E^{(2)}\approx\frac1{15}\bigl(S^{(2)}-S^{(1)}\bigr).
\]
Thus the corrected approximation is
\[
 I\approx S^{(2)}+\frac{S^{(2)}-S^{(1)}}{15}
 =\frac{16S^{(2)}-S^{(1)}}{15}.
\]
If
\[
 \frac{|S^{(2)}-S^{(1)}|}{15}<\varepsilon,
\]
the interval is accepted.  Otherwise it is divided into two halves, each with tolerance $\varepsilon/2$, and the test is repeated recursively.
